\section{Discussion}
\label{sec:discussion}

\subsection{What the multiplier says and does not say}

A multiplier of \mAdv{} on the Adverse base says that, in this model of this bank, the daily run-off rates of the supervisory scenario would have to be about twice as fast for the survival period to fall from \spAdv{} days to 30. It says that the distance to the target is modest, and that the target is not a remote one. It does not say that a run twice as fast is likely or that the bank has that headroom: the deposit behaviour is synthetic, the mapping of the uninsured balance to a supervisory class is a convention, and the multiplier changes by \jointDevAdv{}\,\% across the joint band of conventions. The reverse stress is a measure of distance within a model, and its value to a validator is that it states the distance in the unit of the scenario rather than in days of survival, which differ in how fast they change: \cref{fig:sp-curve} shows that the survival period drops from 62 to 32 days between $m=1$ and $m=2$ and from 13 to 1 day between $m=5$ and $m=21$.

The sensitivity formula explains part of the stability. With an elasticity of \elastAdv{} the multiplier moves by about one and a half times the relative change in capacity, so a shift of 0.05 in every haircut applied to the securities, which changes the capacity by a similar share, moves the multiplier by \hcLowAdvPct{}\,\% or \hcHighAdvPct{}\,\%. The refusals are the more important sensitivity. The slow monetisation schedule and the mapping to other legal entities make the base reach 30 days at a multiplier of 1, and the solver says so instead of returning a multiplier of 1. A table of multipliers across conventions with a column of refusals is more informative than one without, because a refusal marks the convention at which the bank is already at the target.

\subsection{The labelled run}

The score of the fast-run base is red, and the result is not a calibration accident. The maximum implied rate comes from a single row, Northern Rock, and the median and the mean rules give lower ratios. The deeper reason is structural. The base is a constant daily rate on one segment calibrated on runs of 4 to 75 days, and the March 2023 run was a two-day event in which a concentrated uninsured deposit base left at an implied rate of 0.6127 per day. A constant rate fitted on slower runs cannot reproduce a run in which the speed was set by coordination among large depositors, which the empirical work on the run describes \citep{cipriani2024,iyerpuri2012}. The multiplier of the labelled run, \mSvb{} times the calibrated rate against a cap of \capFast{}, quantifies it: the observed run was in the model about \ratioLab{} times further than the multiplier that brings the survival period to at most two days.

What the paper does not claim is as important. The red light does not show that a better calibration rule exists, and the multiplier of the labelled run is not a forecast of the next run. The protocol shows how a stress family can be scored against an event without being calibrated on it, and it records that this family does not reach the event.

\subsection{Failure cases and negative results}

Three predictions failed. The hyperbolic rule of H6 holds for one segment (\cref{prop:hyperbolic}) and fails on a mixture, with a median error of \kHSixMed{}\,\% on the corpus; it is a quick check for a bank dominated by one segment and not a replacement for the solver. The joint stability of H8 failed because of the monetisation schedule: it is the schedule, not the haircut, that moves the base survival period under the target. At the central and high haircuts the slow schedule makes the Adverse base reach the target at a multiplier of 1 and the solver refuses; at the low haircut the base stays above the target, the multiplier is \lowSlowM{}, and the survival period at the multiplier is \lowSlowSp{} days against \lowSlowBelow{} one decimal below it, so the threshold is a cliff and not an undefined value. The whole deviation of \jointDevAdv{}\,\% comes from that corner. The amber bar of H10 failed because of the speed of the run. A fourth result was not predicted and bears on the use of the method: with the monetisation schedule of the case design the survival period of the fast-run base steps from 4 days to 1 across one decimal of the multiplier, so targets of 1, 2 and 3 days have the same multiplier. A reverse stress on a short ladder with a fast-releasing capacity is a statement about a step, and the survival period at the multiplier is the figure that tells it.

\subsection{Implications for validation practice}

A validator who receives a stress model with a table of survival periods under three scenarios can add a column that the paper's method makes available: the multiplier at which each scenario reaches the target the institution cares about, or the statement that it does not. The column does not replace the scenarios, it measures how far they are from the target in one unit. Three practices follow. A reverse stress result should be reported with the survival period at the multiplier and just below it, so that the reader can check the threshold. A refusal should be reported as a result with its reason, in the same table as the multipliers. And a comparison of stress scenarios by their survival periods should count a bank that survives the horizon as longer than any day, as the supervisory exercise does by reporting that the position stays positive; a scan over a coarse grid that returns the first passing point answers the refused instances with a number and, in the corpus, gives a median error of \kGErrPc{}\,\% on the instances it solves.
